\chapter{Assignment 4: Posets and Equivalence Relations}\label{ch:sol4}

\subsection*{Problem 1}
$P=\{1,2,3,4,6,8,12,24\}$ under divisibility (Hasse diagram in Chapter~\ref{ch:orders}); $S=\{2,3,4,6\}$.

\emph{(a) Minimal/maximal in $S$.} Check each $s\in S$ against Definition~\ref{def:extremal}. $2$: the only $s'\in S$ with $s'\mid2$ is $2$ itself — \textbf{minimal}. $3$: likewise only $3\mid3$ in $S$ — \textbf{minimal}. $4$: $2\mid4$ and $2\ne4$, so $4$ is not minimal; is there $s'\in S$ with $4\mid s'$, $s'\ne4$? No ($4\nmid3,6$; wait $4\nmid6$) — so $4$ is \textbf{maximal}. $6$: $2\mid6$ and $3\mid6$, so not minimal; no $s'\in S\setminus\{6\}$ has $6\mid s'$ — \textbf{maximal}.
\[
\text{Minimal}(S)=\{2,3\},\qquad\text{Maximal}(S)=\{4,6\}.
\]

\emph{(b) Minimum/maximum.} Two distinct minimal elements $\Rightarrow$ no minimum (Theorem~\ref{thm:minunique} — a minimum, if it existed, would be the \emph{unique} minimal element). Similarly two distinct maximal elements $\Rightarrow$ \textbf{no maximum}.

\emph{(c) Bounds in $P$.} Lower bounds: $x\in P$ with $x\mid2,x\mid3,x\mid4,x\mid6$, i.e.\ $x\mid\gcd(2,3,4,6)=1$: only $x=1$. \textbf{Lower bounds $=\{1\}$.} Upper bounds: $x\in P$ with $2,3,4,6$ all dividing $x$, i.e.\ $\operatorname{lcm}(2,3,4,6)=12$ divides $x$; among $P$, that means $x\in\{12,24\}$. \textbf{Upper bounds $=\{12,24\}$.}

\emph{(d)} $\glb(S)=\max\{1\}=1$. $\lub(S)=\min\{12,24\}=12$ (since $12\mid24$).

\subsection*{Problem 2}
$X=\{1,2,3,4\}$, $(\pset(X),\subseteq)$, $S=\{\{1\},\{2\},\{1,2\},\{1,3\},\{2,3\}\}$.

\emph{(a)} $\{1\}$: no other member of $S$ is $\subseteq\{1\}$ — \textbf{minimal}. $\{2\}$: likewise \textbf{minimal}. $\{1,2\}$: contains $\{1\}\in S$ properly, not minimal; nothing in $S$ properly contains it — \textbf{maximal}. $\{1,3\}$: contains $\{1\}\in S$ properly, not minimal; nothing in $S$ properly contains $\{1,3\}$ (neither $\{1,2\}$ nor $\{2,3\}$ does) — \textbf{maximal}. $\{2,3\}$: contains $\{2\}$, not minimal; nothing properly contains it — \textbf{maximal}.
\[
\text{Minimal}(S)=\{\{1\},\{2\}\},\qquad\text{Maximal}(S)=\{\{1,2\},\{1,3\},\{2,3\}\}.
\]

\emph{(b)} Two distinct minimal (resp.\ maximal) elements $\Rightarrow$ \textbf{no minimum and no maximum}.

\emph{(c)} Lower bounds: $L\in\pset(X)$ with $L\subseteq$ every member of $S$, i.e.\ $L\subseteq\{1\}\cap\{2\}\cap\{1,2\}\cap\{1,3\}\cap\{2,3\}=\varnothing$: only $L=\varnothing$. \textbf{Lower bounds $=\{\varnothing\}$.} Upper bounds: $U\supseteq\{1\}\cup\{2\}\cup\{1,2\}\cup\{1,3\}\cup\{2,3\}=\{1,2,3\}$: $U\in\{\{1,2,3\},X\}$. \textbf{Upper bounds $=\{\{1,2,3\},X\}$.}

\emph{(d)} $\glb(S)=\varnothing$. $\lub(S)=\{1,2,3\}$ (as $\{1,2,3\}\subseteq X$).

\subsection*{Problem 3}
$\Z\times\Z$, product order (Proposition~\ref{prop:productorder}); $S=\{(0,1),(1,0),(1,2),(2,1)\}$.

Comparabilities: $(0,1)\le(1,2)$ ($0\le1,1\le2$); $(0,1)\le(2,1)$ ($0\le2,1\le1$); $(1,0)\le(1,2)$; $(1,0)\le(2,1)$ ($1\le2,0\le1$). The pairs $(0,1),(1,0)$ are incomparable ($0\le1$ but $1\not\le0$, and vice versa fails too), and $(1,2),(2,1)$ are incomparable.

\emph{(a)} $(0,1)$: nothing in $S$ is $\le$ it (checked above) — \textbf{minimal}. $(1,0)$: likewise \textbf{minimal}. $(1,2)$: $(0,1)\le(1,2)$, so not minimal; nothing in $S$ exceeds it — \textbf{maximal}. $(2,1)$: $(1,0)\le(2,1)$, not minimal; nothing exceeds it — \textbf{maximal}.
\[
\text{Minimal}=\{(0,1),(1,0)\},\qquad\text{Maximal}=\{(1,2),(2,1)\}.
\]

\emph{(b)} Two incomparable minimal (resp.\ maximal) elements $\Rightarrow$ \textbf{no minimum, no maximum}.

\emph{(c)} Lower bounds: $(a,b)$ with $a\le0,1,1,2$ (i.e.\ $a\le0$) and $b\le1,0,2,1$ (i.e.\ $b\le0$): $\{(a,b)\in\Z^2:a\le0,b\le0\}$. Upper bounds: $a\ge\max\{0,1,1,2\}=2$, $b\ge\max\{1,0,2,1\}=2$: $\{(a,b)\in\Z^2:a\ge2,b\ge2\}$.

\emph{(d)} $\glb(S)=(0,0)$ (the maximum of the lower-bound set, itself lower-bound-shaped, attained at $a=0,b=0$). $\lub(S)=(2,2)$.

\subsection*{Problem 4}
$P=\mathrm{Sub}(\R^3)$; $U=\Span\{(1,0,0),(0,1,0)\}$ (the $xy$-plane), $V=\Span\{(0,1,0),(0,0,1)\}$ (the $yz$-plane); $S=\{U,V\}$.

\emph{(a)} $U\ne V$ and neither contains the other ($(1,0,0)\in U\setminus V$; $(0,0,1)\in V\setminus U$), so $U,V$ are incomparable: both \textbf{minimal and maximal} in $S$ (Proposition~\ref{prop:extremal-rel}: with only two, mutually incomparable, elements, each is trivially both).

\emph{(b)} \textbf{No minimum, no maximum} (incomparability).

\emph{(c)} Lower bounds: subspaces $W\subseteq U\cap V$. Now $U\cap V=\Span\{(0,1,0)\}$ (the $y$-axis: a vector in both the $xy$- and $yz$-planes has $x=0$ and $z=0$). The subspaces of a $1$-dimensional space are $\{0\}$ and the space itself, so \textbf{lower bounds $=\{\{0\},\Span\{(0,1,0)\}\}$.} Upper bounds: subspaces $W\supseteq U+V$. Now $U+V=\Span\{(1,0,0),(0,1,0),(0,0,1)\}=\R^3$ (Proposition~\ref{prop:subspaces}), so the only such $W$ is $\R^3$ itself: \textbf{upper bounds $=\{\R^3\}$.}

\emph{(d)} By Proposition~\ref{prop:subspaces}: $\glb(S)=U\cap V=\Span\{(0,1,0)\}$, $\lub(S)=U+V=\R^3$.

\subsection*{Problem 5}
\emph{(a) ``$P\sim Q$ iff $P,Q$ lie on a common line in $\R^2$.''} As discussed in Example~\ref{ex:sameline}: any two points of $\R^2$ lie on a common line (the line through them if distinct, any line through the point if $P=Q$), so $\sim$ is the \emph{total} relation $\R^2\times\R^2$. The total relation on any set is trivially reflexive, symmetric, and transitive, so \textbf{$\sim$ IS an equivalence relation}, with a single equivalence class: $[P]=\R^2$ for every $P$.

\emph{(b) $x\sim y\iff x\mid y$ on $\Z$.} Reflexive: $x\mid x$ (as $x=x\cdot1$). \textbf{Not symmetric}: $2\mid4$ but $4\nmid2$. Hence \textbf{$\sim$ is NOT an equivalence relation.}

\emph{(c) $x\sim y\iff n\mid(x-y)$, fixed $n\in\Z$.} This is Example~\ref{ex:congruence}: reflexive ($n\mid0$), symmetric ($n\mid(x-y)\Rightarrow n\mid(y-x)$), transitive ($n\mid(x-y),n\mid(y-z)\Rightarrow n\mid(x-z)$, as a sum of multiples of $n$). \textbf{$\sim$ IS an equivalence relation.} For $n\ge1$ the classes are $[a]=a+n\Z=\{a+kn:k\in\Z\}$, and by the division algorithm there are exactly $n$ of them: $\Z/n\Z=\{[0],[1],\dots,[n-1]\}$ (for $n=0$, $\sim$ is equality and every class is a singleton; for $n=\pm1$ there is a single class, all of $\Z$).

\subsection*{Problem 6}
This is Example~\ref{ex:refl-trans}: on $A=\{1,2\}$, $R=\{(1,1),(2,2),(1,2)\}$ is reflexive and transitive (both checked exhaustively there) but not symmetric ($(1,2)\in R$, $(2,1)\notin R$).

\subsection*{Problem 7}
This is Example~\ref{ex:trans-only}: on $A=\{1,2\}$, $R=\{(1,2)\}$ is not reflexive ($(1,1)\notin R$), not symmetric ($(2,1)\notin R$), and transitive \emph{vacuously} (no pair $(x,y),(y,z)\in R$ exists with $y$ shared, other than trivial cases, so the implication defining transitivity is never triggered with a false conclusion).

\subsection*{Problem 8}
\emph{(a)} This is Theorem~\ref{thm:minunique}: if $m,m'$ are both minimum elements of a poset, then $m\le m'$ (as $m$ is a minimum) and $m'\le m$ (as $m'$ is), so $m=m'$ by antisymmetry.

\emph{(b)} This is Theorem~\ref{thm:finitemax}, proved there by induction on $|P|$: removing one element, applying the inductive hypothesis to get a maximal element $m'$ of $P\setminus\{p\}$, and checking that either $m'$ or $p$ is maximal in $P$ (analogously for minimal, via the dual order).

\emph{(c)} This is Theorem~\ref{thm:eqdisjoint}, an immediate consequence of Lemma~\ref{lem:classes}: if $[a]\cap[b]\ne\varnothing$ then $[a]=[b]$; so any two classes are equal or (if their intersection is empty) disjoint.

\subsection*{Problem 9}
This is Theorem~\ref{thm:equiv-partition}, proved in full in Chapter~\ref{ch:equiv}:
\begin{itemize}
	\item Given an equivalence relation $\sim$ on $A$, the set of equivalence classes $A/{\sim}$ is a partition: classes are non-empty ($a\in[a]$), pairwise equal-or-disjoint (part (c) above) hence pairwise disjoint when distinct, and their union is $A$ (every $a$ lies in $[a]$).
	\item Given a partition $\mathcal P$ of $A$, the relation $a\sim_{\mathcal P}b:\iff$ ``$a,b$ lie in a common block'' is an equivalence relation, whose classes are exactly the blocks of $\mathcal P$.
	\item The two constructions are mutually inverse, so equivalence relations on $A$ correspond bijectively to partitions of $A$.
\end{itemize}